\slajdid{09_2-s034}
\begin{frame}[label=09_2-s034]
\gusttekst\setlength{\abovedisplayskip}{3pt}\setlength{\belowdisplayskip}{3pt}\setlength{\abovedisplayshortskip}{2pt}\setlength{\belowdisplayshortskip}{2pt}\begin{columns}[c,onlytextwidth]\column{.21\textwidth}\centering\input{slike/tikz/s034_cmos_invertor.tex}\column{.77\textwidth}Česta situacija, jeste $k_n=k_p$ i $V_{Tn}=-V_{Tp}$. U tom slučaju
\[V_I=V_O+\frac{V_{DD}}2\]
pa zamenom u polazni izraz
\[k_p(V_I-V_{DD}-V_{Tp})^2=k_nV_O\left(2(V_I-V_{Tn})-V_O\right)\]
\[\left(V_O+\frac{V_{DD}}2-V_{DD}+V_{Tn}\right)^2=V_O\left(2\left(V_O+\frac{V_{DD}}2-V_{Tn}\right)-V_O\right)\]
\[\left(V_O-\frac{V_{DD}}2+V_{Tn}\right)^2=V_O(V_O+V_{DD}-2V_{Tn})\]
\[V_{O(IH)}=\frac{V_{DD}-2V_{Tn}}8\]
\[V_{IH}=\frac{5V_{DD}-2V_{Tn}}8\]\end{columns}
\end{frame}
